% ---
% id: reductions/wpfss/prop:route
% kind: prop
% title: "Certified weights routed from the receiver artifact"
% summary: >-
%   At absorbing bounded letters the receiver's per-setting high-photon bound
%   survives the dark counts, sharpened by r_min, on every block of the joint
%   family above the cutoff. Rounding it down onto the lattice keeps it
%   certified, so the routed weights are certified weights for the flag-state
%   targets.
% prominence: standard
% status: accepted
% uses: [lem:descent, def:targets, def:weight-route, reductions/def:joint-family, reductions/lem:exactpsd, reductions/def:substrate, prop:weight-monotone, receiver/thm:click-povm, receiver/thm:coarse-graining, receiver/def:slot, receiver/prop:slotvalid, receiver/thm:emitted-bound, receiver/def:setting-family, receiver/cor:unconditioned-family, receiver/detector-imperfections/def:darkadm, receiver/detector-imperfections/thm:dark-counts]
% mentions: [cor:oneabsorbing, req:column-stochastic-refusal, req:descent-decided-exactly, req:fibre-equality, req:absorbing-bounded-letters, foundations/req:rounding-direction]
% assumes: [A-RECV]
% refs: []
% provenance: newly-derived
% lean:
% notes: "Item (iv) says which receiver claim supplies which hypothesis on the emission path, each a uses edge read under A-RECV, and which hypotheses are decided on input instead; its proof is that reading. For a setting family that a receiver submodule constructs, the claims behind it are named by the claim of this submodule that consumes that family."
% ---

The routed weight of Definition~\ref{def:weight-route} is certified for two
reasons: the receiver's per-setting bound survives the dark counts at an
absorbing letter, and the emitted weight is rounded down. Had that rounding
been to nearest, or upwards, the second step would fail: an emitted weight
above $\hat p_{(y_i,t)}$ leaves \eqref{eq:certweights} unproved, and it can
make it false, the failure Assumption~\ref{ass:recv}'s ledger row calls a
security break.
\xref{foundations/req:rounding-direction} names the same direction for every
emitted lower bound. The input conditions part~(iv) leaves to the declaration
are decided by Requirements~\ref{req:column-stochastic-refusal},
\ref{req:descent-decided-exactly}, \ref{req:fibre-equality}
and~\ref{req:absorbing-bounded-letters}.

\begin{proposition}\label{prop:route}
Let $\{M_{y|t}\}_{y\in\Yset}$ be, for each $t\in\Tset$, a POVM on $\Chan$ block
diagonal at the cutoff, and let $R$ descend through $\cg$ with descended matrix
$R^{\cg}$ (Lemma~\ref{lem:descent}). Let $\{q_t\}$ be a probability
distribution, and let $\{(\mathcal{O}_i,(p^{(t)}_i)_{t\in\Tset})\}_{i=1}^{r}$
be the declared slot, with $\mathcal{O}_i=\cg^{-1}(y_i)$ for pairwise distinct
letters $y_1,\dots,y_r$, with $0\le p^{(t)}_i\le1$, and with the per-setting
guarantee $\blk{M}{\gtN}_{y_i|t}\succeq p^{(t)}_i\,\idon{}$ for every
$t\in\Tset$ and every $i$. Suppose each bounded letter is \emph{absorbing},
\begin{equation}\label{eq:absorbing}
  R^{\cg}_{y_i|y_i}=1
  \qquad\text{for every }i=1,\dots,r,
  \tag{absorbing}
\end{equation}
and put $r_{\min,i}\coloneqq\min_{y'\ne y_i}R^{\cg}_{y_i|y'}$, with the
convention $r_{\min,i}\coloneqq0$ when the minimum is over the empty set (that
is, when $\Yset=\{y_i\}$), the conservative value, which reduces
\eqref{eq:routebound} to the un-improved bound. Then:
\begin{enumerate}[label=(\roman*),leftmargin=2.4em]
\item The joint-alphabet POVM
  $M_{(y,t)}\coloneqq q_t\sum_{y'}R^{\cg}_{y|y'}M_{y'|t}$ of
  Definitions~\ref{def:joint-family} and~\ref{def:weight-route} satisfies, for
  every $i$ and every $t\in\Tset$,
  \begin{equation}\label{eq:routebound}
    \blk{M}{\gtN}_{(y_i,t)}\ \succeq\
    q_t\Big(\big(1-r_{\min,i}\big)\,p^{(t)}_i+r_{\min,i}\Big)\,\idon{}
    \;=\;\hat p_{(y_i,t)}\,\idon{} .
  \end{equation}
  A slot whose weights do not depend on the setting, $p^{(t)}_i=p_i$ for
  every $t$ (the $t$-uniform form, which every one-setting slot has), is the
  special case, and there \eqref{eq:routebound} is the bound in $p_i$.
\item The \emph{emitted} weight of \eqref{eq:weightroute} is certified, because
  the rounding runs downwards. Every number on the right of
  \eqref{eq:routebound} lies on the declared lattice: $q_t$ and $p^{(t)}_i$ by
  declaration, and $r_{\min,i}$ as an entry of $R^{\cg}$, which is a sum of
  declared entries and lies in $[0,1]$. So $\hat p_{(y_i,t)}$ is an exact
  element of $\Fld$ formed from them; and
  $p_{(y_i,t)}=\lat{\hat p_{(y_i,t)}}\le\hat p_{(y_i,t)}$ by
  \eqref{eq:latdown}, whence
  \begin{equation*}
    \blk{M}{\gtN}_{(y_i,t)}\ \succeq\ p_{(y_i,t)}\,\idon{} ,
  \end{equation*}
  which is \eqref{eq:certroute}; so the emitted $(p_v)$ are
  \tref{certified weights} in the sense of Definition~\ref{def:targets} and they lie on the declared
  lattice.
\item Under the sufficient condition \eqref{eq:canonicalcg} of
  Lemma~\ref{lem:descent} and with the single entry $\mathcal{O}_1=\MCK$,
  \eqref{eq:absorbing} holds at $y_1=y_{\mc}$, and its per-setting hypothesis
  is the supplier's slot guarantee
  $\sum_{S\in\MCK}\blk{M}{\gtN}_{S|t}\succeq p^{(t)}_1\,\idon{}$ at each $t$,
  read through $\cg^{-1}(y_{\mc})=\MCK$.
\item On the emission path the hypotheses about the supplier's objects hold
  under Assumption~\ref{ass:recv}: the conditional families, and their
  block diagonality in the photon number, block by block and so in
  particular at the cutoff, are at every setting those of the coarse-grained
  elements of the \tref{setting family} on which the supplier states its
  weights, by item~(a) of \xref{receiver/def:setting-family} and
  \xref{receiver/thm:coarse-graining}; for a receiver without probe inputs
  that family is the one-setting family of
  \xref{receiver/cor:unconditioned-family}, and the conditional families and
  their block diagonality are given by the receiver's
  \xref{receiver/thm:click-povm}\,(iii) and
  \xref{receiver/thm:coarse-graining}; the
  dark-count matrix $R$, and the block diagonality in the photon number of
  the elements with the dark counts composed, by
  \xref{receiver/detector-imperfections/def:darkadm} and
  \xref{receiver/detector-imperfections/thm:dark-counts}; so every element of
  the joint family of (i), and of the post-processed joint family that
  Definition~\ref{def:weight-route} forms in either of its forms, is block
  diagonal in the photon number in every block, below the cutoff as above
  it; and the slot, its
  per-setting guarantee (the first conclusion of
  \xref{receiver/thm:emitted-bound}) and $0\le p^{(t)}_i\le1$,
  $\sum_ip^{(t)}_i\le1$ at every $t$ by \xref{receiver/def:slot},
  \xref{receiver/thm:emitted-bound} and \xref{receiver/prop:slotvalid}. The
  other hypotheses are not supplier facts: that $R$ is column-stochastic and descends through $\cg$, that
  $\mathcal{O}_i=\cg^{-1}(y_i)$ and \eqref{eq:absorbing} are decided exactly
  on the declared data, a declaration failing one being
  refused, and $\{q_t\}$ are the declared \tref{class weights}
  (Definition~\ref{def:joint-family}).
\end{enumerate}
\end{proposition}

\begin{proof}
(i) Fix $i$ and $t$. Split the sum over columns and use \eqref{eq:absorbing}
at that letter:
\begin{align*}
  \sum_{y'}R^{\cg}_{y_i|y'}M_{y'|t}
  &=M_{y_i|t}+\sum_{y'\ne y_i}R^{\cg}_{y_i|y'}M_{y'|t}
  \;\succeq\;M_{y_i|t}+r_{\min,i}\sum_{y'\ne y_i}M_{y'|t}\\
  &=(1-r_{\min,i})\,M_{y_i|t}+r_{\min,i}\,\idon{},
\end{align*}
using $M_{y'|t}\succeq0$ and $R^{\cg}_{y_i|y'}\ge r_{\min,i}$ in the
inequality and POVM completeness in the last step. All operators are block
diagonal at the cutoff, so the inequality restricts to the $\gtN$ block, where
$\blk{M}{\gtN}_{y_i|t}\succeq p^{(t)}_i\,\idon{}$ and $1-r_{\min,i}\ge0$, every
entry of $R^{\cg}$ lying in $[0,1]$ by non-negativity and unit columns
(Lemma~\ref{lem:descent}). Multiplying by $q_t\ge0$ gives
\eqref{eq:routebound}. Nothing in the argument couples the entries: each $i$
uses only its own row of $R^{\cg}$ and its own weight, which is why $r_{\min}$
must be indexed (Corollary~\ref{cor:oneabsorbing} says what that costs).

(ii) $\hat p_{(y_i,t)}$ is formed from $q_t$, $p^{(t)}_i$ and $r_{\min,i}$ by
finitely many field operations in $\Fld$ (Lemma~\ref{lem:exactpsd},
Definition~\ref{def:substrate}), so it \emph{is} the real value of
\eqref{eq:weightroute}'s expression and \eqref{eq:routebound} is a statement
about that value. Monotonicity of $\succeq$ in a scalar multiple of the
identity, with $p_{(y_i,t)}\le\hat p_{(y_i,t)}$, gives \eqref{eq:certroute};
and $p_{(y_i,t)}\ge0$ because $\hat p_{(y_i,t)}\ge0$ and $\lat{\cdot}$ carries
a non-negative rational to a non-negative lattice value. No rounding-error
budget is tracked and none is needed: the inequality is exact, and the only
inexactness anywhere near it is a rounding that weakens it
(Proposition~\ref{prop:weight-monotone}\,(i)).

(iii) The proof of Lemma~\ref{lem:descent} under \eqref{eq:canonicalcg} gives
$P(y_{\mc}|S')=1$ for every $S'\in\MCK$, hence
$R^{\cg}_{y_{\mc}|y_{\mc}}=1$; and
$M_{y_{\mc}|t}=\sum_{S\in\cg^{-1}(y_{\mc})}M_{S|t}=\sum_{S\in\MCK}M_{S|t}$ by
the definition of the coarse-grained POVM.

(iv) Each receiver claim cited there states the hypothesis it is cited for,
of the objects the receiver artifact supplies, and
Assumption~\ref{ass:recv} makes the artifact's facts hold of the numbers as
declared. The supplier states its weights on a setting family
(\xref{receiver/thm:emitted-bound}), whose elements at every setting commute
with the photon number by item~(a) of \xref{receiver/def:setting-family},
and a coarse-grained element is a sum of them
(\xref{receiver/thm:coarse-graining}); for a receiver without probe inputs
\xref{receiver/cor:unconditioned-family} identifies that family as the
click-pattern POVM of \xref{receiver/thm:click-povm}, block diagonal by its
item~(iii). The remaining hypotheses are the input decisions and the
declaration named there. An element of the joint family of (i), and an
element of the post-processed joint family in either form, is a combination
with non-negative coefficients ($q_t$ and the entries of $R^{\cg}$, or $q_t$
alone) of the receiver's elements, conditional or with the dark counts
composed, which the claims cited give block diagonal in the photon number;
and a sum of operators each block diagonal in the photon number is again
block diagonal in it, block by block.
\end{proof}
